% Oversells «Test» tugmasi yaratgan fayl: savollar bankidan variantlar, javob varaqasi va javoblar kaliti.
% Qoʻlda tahrirlamang: test qayta yaratilganda bu fayl yangidan yoziladi; savollarni bank faylida oʻzgartiring.
\documentclass[11pt,a4paper]{article}
\usepackage{fontspec}
\setmainfont{Noto Serif}
\usepackage{babel}
\babelprovide[import=uz-Latn, main]{uzbek}
\def\uzbekhyphenmins{22}
\usepackage{amsmath,amssymb}
\usepackage[top=1.8cm,bottom=1.8cm,left=2cm,right=2cm]{geometry}
\usepackage{testlar}
\begin{document}

% --- A-variant ---
\testvariant{Matematika. Nazorat ishi}{A}
\testsavol{$x^2-5x+6=0$ tenglamaning ildizlarini toping.}{4}{%
  \testjavob{A}{$-2$ va $-3$}%
  \testjavob{B}{$-1$ va $6$}%
  \testjavob{C}{$1$ va $6$}%
  \testjavob{D}{$2$ va $3$}%
}
\testsavol{Toʻgʻri burchakli uchburchakning katetlari $6$ va $8$ ga teng. Gipotenuzani toping.}{4}{%
  \testjavob{A}{$10$}%
  \testjavob{B}{$\sqrt{28}$}%
  \testjavob{C}{$14$}%
  \testjavob{D}{$12$}%
}
\testsavol{$\log_2 32$ ning qiymati nechaga teng?}{4}{%
  \testjavob{A}{$16$}%
  \testjavob{B}{$4$}%
  \testjavob{C}{$5$}%
  \testjavob{D}{$6$}%
}
\testsavol{Arifmetik progressiyada $a_1=3$, $d=4$. $a_{10}$ ni toping.}{4}{%
  \testjavob{A}{$36$}%
  \testjavob{B}{$40$}%
  \testjavob{C}{$39$}%
  \testjavob{D}{$43$}%
}
\testsavol{$2+3\cdot 4$ ifodaning qiymatini toping.}{4}{%
  \testjavob{A}{$24$}%
  \testjavob{B}{$20$}%
  \testjavob{C}{$9$}%
  \testjavob{D}{$14$}%
}
\testsavol{$f(x)=x^3$ funksiyaning hosilasini toping.}{4}{%
  \testjavob{A}{$\dfrac{x^4}{4}$}%
  \testjavob{B}{$3x^3$}%
  \testjavob{C}{$x^2$}%
  \testjavob{D}{$3x^2$}%
}
\testsavol{Qaysi son tub son?}{4}{%
  \testjavob{A}{$29$}%
  \testjavob{B}{$21$}%
  \testjavob{C}{$27$}%
  \testjavob{D}{$33$}%
}
\testsavol{Doiraning radiusi $5$ sm. Uning yuzini toping.}{4}{%
  \testjavob{A}{$10\pi$ sm$^2$}%
  \testjavob{B}{$5\pi$ sm$^2$}%
  \testjavob{C}{$50\pi$ sm$^2$}%
  \testjavob{D}{$25\pi$ sm$^2$}%
}
\testsavol{$\dfrac{1}{2}+\dfrac{1}{3}$ yigʻindini hisoblang.}{4}{%
  \testjavob{A}{$\dfrac{2}{6}$}%
  \testjavob{B}{$\dfrac{2}{5}$}%
  \testjavob{C}{$\dfrac{5}{6}$}%
  \testjavob{D}{$\dfrac{1}{6}$}%
}
\testsavol{$\sin^2\alpha+\cos^2\alpha$ ifoda nimaga teng?}{4}{%
  \testjavob{A}{$2$}%
  \testjavob{B}{$0$}%
  \testjavob{C}{$\sin 2\alpha$}%
  \testjavob{D}{$1$}%
}

% --- B-variant ---
\testvariant{Matematika. Nazorat ishi}{B}
\testsavol{$\sin^2\alpha+\cos^2\alpha$ ifoda nimaga teng?}{4}{%
  \testjavob{A}{$0$}%
  \testjavob{B}{$\sin 2\alpha$}%
  \testjavob{C}{$1$}%
  \testjavob{D}{$2$}%
}
\testsavol{Arifmetik progressiyada $a_1=3$, $d=4$. $a_{10}$ ni toping.}{4}{%
  \testjavob{A}{$43$}%
  \testjavob{B}{$39$}%
  \testjavob{C}{$36$}%
  \testjavob{D}{$40$}%
}
\testsavol{Qaysi son tub son?}{4}{%
  \testjavob{A}{$33$}%
  \testjavob{B}{$29$}%
  \testjavob{C}{$27$}%
  \testjavob{D}{$21$}%
}
\testsavol{$2+3\cdot 4$ ifodaning qiymatini toping.}{4}{%
  \testjavob{A}{$24$}%
  \testjavob{B}{$9$}%
  \testjavob{C}{$14$}%
  \testjavob{D}{$20$}%
}
\testsavol{$x^2-5x+6=0$ tenglamaning ildizlarini toping.}{4}{%
  \testjavob{A}{$-2$ va $-3$}%
  \testjavob{B}{$1$ va $6$}%
  \testjavob{C}{$2$ va $3$}%
  \testjavob{D}{$-1$ va $6$}%
}
\testsavol{$f(x)=x^3$ funksiyaning hosilasini toping.}{4}{%
  \testjavob{A}{$x^2$}%
  \testjavob{B}{$3x^2$}%
  \testjavob{C}{$\dfrac{x^4}{4}$}%
  \testjavob{D}{$3x^3$}%
}
\testsavol{$\dfrac{1}{2}+\dfrac{1}{3}$ yigʻindini hisoblang.}{4}{%
  \testjavob{A}{$\dfrac{2}{6}$}%
  \testjavob{B}{$\dfrac{1}{6}$}%
  \testjavob{C}{$\dfrac{5}{6}$}%
  \testjavob{D}{$\dfrac{2}{5}$}%
}
\testsavol{$\log_2 32$ ning qiymati nechaga teng?}{4}{%
  \testjavob{A}{$4$}%
  \testjavob{B}{$6$}%
  \testjavob{C}{$16$}%
  \testjavob{D}{$5$}%
}
\testsavol{Doiraning radiusi $5$ sm. Uning yuzini toping.}{4}{%
  \testjavob{A}{$25\pi$ sm$^2$}%
  \testjavob{B}{$50\pi$ sm$^2$}%
  \testjavob{C}{$10\pi$ sm$^2$}%
  \testjavob{D}{$5\pi$ sm$^2$}%
}
\testsavol{Toʻgʻri burchakli uchburchakning katetlari $6$ va $8$ ga teng. Gipotenuzani toping.}{4}{%
  \testjavob{A}{$12$}%
  \testjavob{B}{$14$}%
  \testjavob{C}{$10$}%
  \testjavob{D}{$\sqrt{28}$}%
}

% --- C-variant ---
\testvariant{Matematika. Nazorat ishi}{C}
\testsavol{$\log_2 32$ ning qiymati nechaga teng?}{4}{%
  \testjavob{A}{$16$}%
  \testjavob{B}{$4$}%
  \testjavob{C}{$5$}%
  \testjavob{D}{$6$}%
}
\testsavol{Arifmetik progressiyada $a_1=3$, $d=4$. $a_{10}$ ni toping.}{4}{%
  \testjavob{A}{$43$}%
  \testjavob{B}{$40$}%
  \testjavob{C}{$36$}%
  \testjavob{D}{$39$}%
}
\testsavol{$\dfrac{1}{2}+\dfrac{1}{3}$ yigʻindini hisoblang.}{4}{%
  \testjavob{A}{$\dfrac{2}{6}$}%
  \testjavob{B}{$\dfrac{1}{6}$}%
  \testjavob{C}{$\dfrac{5}{6}$}%
  \testjavob{D}{$\dfrac{2}{5}$}%
}
\testsavol{Toʻgʻri burchakli uchburchakning katetlari $6$ va $8$ ga teng. Gipotenuzani toping.}{4}{%
  \testjavob{A}{$14$}%
  \testjavob{B}{$10$}%
  \testjavob{C}{$\sqrt{28}$}%
  \testjavob{D}{$12$}%
}
\testsavol{Qaysi son tub son?}{4}{%
  \testjavob{A}{$27$}%
  \testjavob{B}{$29$}%
  \testjavob{C}{$33$}%
  \testjavob{D}{$21$}%
}
\testsavol{$f(x)=x^3$ funksiyaning hosilasini toping.}{4}{%
  \testjavob{A}{$x^2$}%
  \testjavob{B}{$\dfrac{x^4}{4}$}%
  \testjavob{C}{$3x^2$}%
  \testjavob{D}{$3x^3$}%
}
\testsavol{$x^2-5x+6=0$ tenglamaning ildizlarini toping.}{4}{%
  \testjavob{A}{$-2$ va $-3$}%
  \testjavob{B}{$-1$ va $6$}%
  \testjavob{C}{$1$ va $6$}%
  \testjavob{D}{$2$ va $3$}%
}
\testsavol{Doiraning radiusi $5$ sm. Uning yuzini toping.}{4}{%
  \testjavob{A}{$5\pi$ sm$^2$}%
  \testjavob{B}{$25\pi$ sm$^2$}%
  \testjavob{C}{$50\pi$ sm$^2$}%
  \testjavob{D}{$10\pi$ sm$^2$}%
}
\testsavol{$\sin^2\alpha+\cos^2\alpha$ ifoda nimaga teng?}{4}{%
  \testjavob{A}{$0$}%
  \testjavob{B}{$\sin 2\alpha$}%
  \testjavob{C}{$2$}%
  \testjavob{D}{$1$}%
}
\testsavol{$2+3\cdot 4$ ifodaning qiymatini toping.}{4}{%
  \testjavob{A}{$24$}%
  \testjavob{B}{$9$}%
  \testjavob{C}{$20$}%
  \testjavob{D}{$14$}%
}

% --- D-variant ---
\testvariant{Matematika. Nazorat ishi}{D}
\testsavol{$\log_2 32$ ning qiymati nechaga teng?}{4}{%
  \testjavob{A}{$4$}%
  \testjavob{B}{$6$}%
  \testjavob{C}{$5$}%
  \testjavob{D}{$16$}%
}
\testsavol{Doiraning radiusi $5$ sm. Uning yuzini toping.}{4}{%
  \testjavob{A}{$25\pi$ sm$^2$}%
  \testjavob{B}{$10\pi$ sm$^2$}%
  \testjavob{C}{$50\pi$ sm$^2$}%
  \testjavob{D}{$5\pi$ sm$^2$}%
}
\testsavol{Toʻgʻri burchakli uchburchakning katetlari $6$ va $8$ ga teng. Gipotenuzani toping.}{4}{%
  \testjavob{A}{$10$}%
  \testjavob{B}{$\sqrt{28}$}%
  \testjavob{C}{$12$}%
  \testjavob{D}{$14$}%
}
\testsavol{$x^2-5x+6=0$ tenglamaning ildizlarini toping.}{4}{%
  \testjavob{A}{$2$ va $3$}%
  \testjavob{B}{$1$ va $6$}%
  \testjavob{C}{$-2$ va $-3$}%
  \testjavob{D}{$-1$ va $6$}%
}
\testsavol{Arifmetik progressiyada $a_1=3$, $d=4$. $a_{10}$ ni toping.}{4}{%
  \testjavob{A}{$40$}%
  \testjavob{B}{$36$}%
  \testjavob{C}{$43$}%
  \testjavob{D}{$39$}%
}
\testsavol{$2+3\cdot 4$ ifodaning qiymatini toping.}{4}{%
  \testjavob{A}{$24$}%
  \testjavob{B}{$14$}%
  \testjavob{C}{$9$}%
  \testjavob{D}{$20$}%
}
\testsavol{Qaysi son tub son?}{4}{%
  \testjavob{A}{$21$}%
  \testjavob{B}{$27$}%
  \testjavob{C}{$33$}%
  \testjavob{D}{$29$}%
}
\testsavol{$\dfrac{1}{2}+\dfrac{1}{3}$ yigʻindini hisoblang.}{4}{%
  \testjavob{A}{$\dfrac{2}{5}$}%
  \testjavob{B}{$\dfrac{1}{6}$}%
  \testjavob{C}{$\dfrac{5}{6}$}%
  \testjavob{D}{$\dfrac{2}{6}$}%
}
\testsavol{$\sin^2\alpha+\cos^2\alpha$ ifoda nimaga teng?}{4}{%
  \testjavob{A}{$2$}%
  \testjavob{B}{$0$}%
  \testjavob{C}{$\sin 2\alpha$}%
  \testjavob{D}{$1$}%
}
\testsavol{$f(x)=x^3$ funksiyaning hosilasini toping.}{4}{%
  \testjavob{A}{$x^2$}%
  \testjavob{B}{$3x^3$}%
  \testjavob{C}{$\dfrac{x^4}{4}$}%
  \testjavob{D}{$3x^2$}%
}

% --- Javob varaqasi ---
\testjavobvaraqasi{Matematika. Nazorat ishi}{%
\begin{tikzpicture}[x=0.85cm,y=-0.85cm,every node/.style={font=\small}]
\node[anchor=west,font=\small\bfseries] at (0,0) {Variant:};
\testpufak{3.2}{0}{A}
\testpufak{4.2}{0}{B}
\testpufak{5.2}{0}{C}
\testpufak{6.2}{0}{D}
\node[anchor=east,font=\small\bfseries] at (1.0,2) {1.};
\testpufak{1.6}{2}{A}
\testpufak{2.6}{2}{B}
\testpufak{3.6}{2}{C}
\testpufak{4.6}{2}{D}
\node[anchor=east,font=\small\bfseries] at (1.0,3) {2.};
\testpufak{1.6}{3}{A}
\testpufak{2.6}{3}{B}
\testpufak{3.6}{3}{C}
\testpufak{4.6}{3}{D}
\node[anchor=east,font=\small\bfseries] at (1.0,4) {3.};
\testpufak{1.6}{4}{A}
\testpufak{2.6}{4}{B}
\testpufak{3.6}{4}{C}
\testpufak{4.6}{4}{D}
\node[anchor=east,font=\small\bfseries] at (1.0,5) {4.};
\testpufak{1.6}{5}{A}
\testpufak{2.6}{5}{B}
\testpufak{3.6}{5}{C}
\testpufak{4.6}{5}{D}
\node[anchor=east,font=\small\bfseries] at (1.0,6) {5.};
\testpufak{1.6}{6}{A}
\testpufak{2.6}{6}{B}
\testpufak{3.6}{6}{C}
\testpufak{4.6}{6}{D}
\node[anchor=east,font=\small\bfseries] at (1.0,7) {6.};
\testpufak{1.6}{7}{A}
\testpufak{2.6}{7}{B}
\testpufak{3.6}{7}{C}
\testpufak{4.6}{7}{D}
\node[anchor=east,font=\small\bfseries] at (1.0,8) {7.};
\testpufak{1.6}{8}{A}
\testpufak{2.6}{8}{B}
\testpufak{3.6}{8}{C}
\testpufak{4.6}{8}{D}
\node[anchor=east,font=\small\bfseries] at (1.0,9) {8.};
\testpufak{1.6}{9}{A}
\testpufak{2.6}{9}{B}
\testpufak{3.6}{9}{C}
\testpufak{4.6}{9}{D}
\node[anchor=east,font=\small\bfseries] at (1.0,10) {9.};
\testpufak{1.6}{10}{A}
\testpufak{2.6}{10}{B}
\testpufak{3.6}{10}{C}
\testpufak{4.6}{10}{D}
\node[anchor=east,font=\small\bfseries] at (1.0,11) {10.};
\testpufak{1.6}{11}{A}
\testpufak{2.6}{11}{B}
\testpufak{3.6}{11}{C}
\testpufak{4.6}{11}{D}
\end{tikzpicture}
}

% --- Javoblar kaliti ---
\testkalit{Matematika. Nazorat ishi}{%
\begin{tabular}{|c|c|c|c|c|c|c|c|c|c|c|}\hline
\textbf{Variant} & \textbf{1} & \textbf{2} & \textbf{3} & \textbf{4} & \textbf{5} & \textbf{6} & \textbf{7} & \textbf{8} & \textbf{9} & \textbf{10} \\\hline
A & D & A & C & C & D & D & A & D & C & D \\\hline
B & C & B & B & C & C & B & C & D & A & C \\\hline
C & C & D & C & B & B & C & D & B & D & D \\\hline
D & C & A & A & A & D & B & D & C & D & D \\\hline
\end{tabular}\par\bigskip
}

\end{document}
